\section*{Solutions to Exercises From Section 17}
\begin{solution}\label{sol:2.17.1}\ref{ex:2.17.1}
    \begin{enumerate}
        \item Every squared real number is greater than $0$. This is false, since $0^{2}=0$.
        \item Any real number, raised to the power of some natural number is less than or equal to $0$. This statement is false. Consider $x=8,n=2$, then $x^{n}\ge 0$. 
        \item There is a real number such that for any other real number, $(ax=x)$. This is true, consider $a=1$.
        \item For any set in the powerset of the natural numbers, that set is a subset of the real numbers. This is also true, since $\N\subseteq\R$. 
        \item For any natural number, there is a member of the powerset of natural numbers such that the cardinality of that set is less than that natural number. This is false, consider a set pick $n=0$, then if $\lvert X \rvert=0$ the statement would not hold. If we suppose that $0\notin\N$, then the statement would be true. We can pick the lowest value for $n=1$, and for the cardinality of $X$, and $0<1$. We can pick $X=\varnothing$ and $\lvert X \rvert=0$, since we are looking at subsets of the $\powerset(\N)$.  
        \item There is a natural number such for any member of the powerset of natural numbers, the cardinality of the set is less than that natural number. This is false. There is no one fixed natural number which is strictly greater than the cardinality of all subsets of the natural numbers. $\powerset(\N)$ contains infinitely many subsets, so it has cardinality $2^{\aleph_{0}}$, whereas $\lvert \N \rvert=\aleph_0$ so we cannot claim that $\aleph_{0}<n$ for some natural number. 
        \item For any subset of the natural numbers there is some integer such that the cardinality of the set is equivalent to that integer. This is also false, the reasoning is simialr to the above example. 
        \item For any integer, there is a subset of subset of the natural numbers, such that the cardinality of that subset is equivalent to an arbitrary integer. This is false, consider $n=-1$. 
        \item For any integer, there is some other integer such that it is the sum of the first integer and $5$. This is true. Consider an arbitrary integer $n=4$, then $m=4+5$ and $m=9$. 
        \item There is some integer, such that for any integer we want, we can find the sum of the first by adding $5$ to the second integer. This is false. This is backwards, suppose we pick some integer $m$, then the statement claims that any integer we insert in the spot of $n$ should give us $m=n+5$. 
    \end{enumerate}
\end{solution}